\chapter{Geometry}\label{ch:geometry}

\section{Requirements}\label{sec:geometry-requirements}

\begin{longtable}{@{}P{0.16\linewidth} P{0.36\linewidth} P{0.40\linewidth}@{}}
\caption{Requirements for the STL file.}\label{tab:stl-req}\\
\toprule
Item & Requirement & Reason \\
\midrule
\endfirsthead
\toprule
Item & Requirement & Reason \\
\midrule
\endhead
Format & ASCII STL & A binary STL has no solid names. All the rocket then becomes one patch. \\
Solid names & \val{nosecone}, \val{body}, \val{boattail}, \val{fins} & \file{Allmesh} changes the first three to \val{cone}, \val{walls}, \val{tail}. The dictionaries and the cases use these four names. \\
Units & Metres & \file{meshDict} and \file{meshInfo} use metres. \\
Position & Nose tip at the origin, axis along $+x$ & The refinement regions, the fin region in \file{meshInfo} and the moment reference \ent{CofR (0 0 0)} use this position. \\
Edge files & \file{<name>_finEdges.obj} and \file{<name>_finRoots.obj} next to the STL & \file{Allmesh} copies them for \ent{edgeMeshRefinement} (\cref{sec:edge-refinement}). \\
\bottomrule
\end{longtable}

\subsection{Solid names that are not in the list}

\file{Allmesh} does not stop if a solid has a different name. The patch keeps the
name of the solid. \Cref{tab:bad-names} gives the effects.

\begin{longtable}{@{}P{0.32\linewidth} P{0.62\linewidth}@{}}
\caption{Effects of a solid name that is not in the list.}\label{tab:bad-names}\\
\toprule
Step & Effect \\
\midrule
\endfirsthead
\toprule
Step & Effect \\
\midrule
\endhead
\cmd{cartesianMesh}, \ent{localRefinement} & An entry without a patch gives only a warning. The patch gets no surface refinement. The cell size near the wall is then that of the volume region (\unit{50}{mm} in \ent{nearField}), not \unit{6.25}{mm}. \\
\cmd{cartesianMesh}, \ent{boundaryLayers} & The patch gets only the global layers (\cref{sec:layers}). \\
Case, \file{0.orig/} & The boundary conditions use the names \val{cone}, \val{walls}, \val{tail}, \val{fins}. A patch with a different name has no boundary condition, and the solver stops. \\
Case, \ent{forceCoeffs} & \ent{wallPatches} in \file{meshInfo} does not contain the patch. Its force is not in the coefficients. \\
\bottomrule
\end{longtable}

Before you make a mesh from a new STL file, examine the solid names:

\begin{shell}
grep -E '^ *solid ' mesh/stl/<file>.stl
\end{shell}

\section{makeRocket.py}\label{sec:makerocket}

\file{mesh/makeRocket.py} writes the geometry from its dimensions. It writes three
files: the STL file, \file{<name>_finEdges.obj} (the leading and trailing edges of
the fins) and \file{<name>_finRoots.obj} (the root of each fin). \cmd{python3
mesh/makeRocket.py -h} gives the options. A file with one \texttt{name value} on each
line can give the dimensions (\ent{--dims}).

To write the geometry of the full mesh (fins at $0^\circ$):

\begin{shell}
python3 mesh/makeRocket.py mesh/stl/Aconcagua_biconvex117.stl
\end{shell}

To write the geometry of the half mesh (fins at $45^\circ$, the default of
\file{Allmesh}):

\begin{shell}
python3 mesh/makeRocket.py mesh/stl/Aconcagua_biconvex117_roll45.stl --finRoll 45
\end{shell}

\ent{--finRoll 45} turns the fins to $45^\circ$, $135^\circ$, $225^\circ$ and
$315^\circ$. The plane $z = 0$ then passes between two fins. The STL file is always
the whole rocket: the half mesh is cut from the full mesh (\cref{sec:half-mesh}).

\begin{table}[htbp]
\centering
\caption{Parts of the Aconcagua geometry.}\label{tab:parts}
\begin{tabularx}{\linewidth}{@{}l L l@{}}
\toprule
Part & Shape & Solid name \\
\midrule
Nose & Von Kármán, \unit{0.8}{m} long & \val{nosecone} \\
Body & Cylinder, $d = \unit{0.151}{m}$ & \val{body} \\
Fins & 4 fins, biconvex section of constant thickness \unit{11.7}{mm}, edges rounded with \unit{1}{mm}, a fillet of \unit{3}{mm} all round the root & \val{fins} \\
Boattail & Cone to $d = \unit{0.110}{m}$, and the base. Shoulder and base edge rounded with \unit{3}{mm} & \val{boattail} \\
\bottomrule
\end{tabularx}
\end{table}

\subsection{Properties of the surface}

The script writes a file only if the surface is closed and all the normals point
out. The properties in \cref{tab:surface-props} come from measurements with
\cmd{cartesianMesh}. Do not change them.

\begin{longtable}{@{}P{0.40\linewidth} P{0.54\linewidth}@{}}
\caption{Properties of the surface and the measurements that set them.}\label{tab:surface-props}\\
\toprule
Property & Reason, measured with the same \file{meshDict} \\
\midrule
\endfirsthead
\toprule
Property & Reason \\
\midrule
\endhead
Nose and body + fins + boattail are two closed shells that touch on coplanar caps. The rings of the two shells are turned by half a facet. & With a sharp shoulder, one continuous hull gave 12 to 87 negative cell volumes at the shoulder. With the shoulder rounded (\ent{--shoulderRadius}, \unit{3}{mm}) the continuous hull gave 0 on $h1$. The boattail must be in the body shell, because the root fillet of the trailing edge runs \unit{3}{mm} onto the cone. \\
Body and fins are one shell: the fin root is cut into the cylinder. & With the fins as separate shells sunk into the body, the fin-body junction was not an edge of the surface. The layers of the two surfaces met at one node in the corner. \\
The leading edge, the trailing edge and the tip of the fins have a radius (\ent{--finEdgeRadius}, \unit{1}{mm}). & With sharp edges, the layers of the two fin faces met at the edge: maximum skewness 25 to 33, non-orthogonality up to $176^\circ$. With the radius: skewness 3.7, non-orthogonality $77^\circ$ to $82^\circ$. \\
A fillet of \unit{3}{mm} all round each fin root (\ent{--finRootRadius}). & Without it, the fin face and the body met at about $90^\circ$, and the fin layers and the body layers closed to one node along the root. With it, on $h1$: non-orthogonality $81.9^\circ \to 67.7^\circ$, skewness $3.58 \to 1.89$, smallest cell $2.6 \to 6.9 \times 10^{-13}$\,m\textsuperscript{3}. \\
The base edge is rounded with \unit{3}{mm} (\ent{--baseEdgeRadius}). & \\
The root trailing edge is on the shoulder (\ent{--finFromBase} = \ent{--tailLength}). & \\
44 facets around the axis (\ent{--nTheta}). & With 96 facets the continuous hull had more negative volumes. \\
\bottomrule
\end{longtable}

The edge radius is in the section parallel to the axis. Normal to the swept leading
edge ($57^\circ$), it is approximately \unit{0.5}{mm}. The section is the convex hull
of the biconvex arcs and a circle of the radius at each edge. The profile is thus a
little thicker in approximately the first \unit{80}{mm} of the chord.

\section{Use a new geometry}

\begin{steps}
\item Write the STL file with \file{makeRocket.py} (\cref{sec:makerocket}), or put
  it in \file{mesh/stl/} with its two edge files.
\item Make sure that the STL file agrees with \cref{tab:stl-req}.
\item Change the six distances of \cmd{surfaceGenerateBoundingBox} in
  \file{Allmesh} (\cref{sec:domain}).
\item Write \ent{lBody} and \ent{rBody} of the new rocket in \file{mesh/meshInfo}.
  The refinement regions scale with them (\cref{sec:object-refinements}).
\item Change the other entries of \file{mesh/meshInfo} (\cref{sec:meshinfo}).
\item Make the mesh: \cmd{mesh/Allmesh <mesh-dir> mesh/stl/<file>.stl}.
\item Remove the old meshes. \file{runMesh.sh} does not compare the geometry
  (\cref{sec:run-mesh}).
\end{steps}
